\documentclass[10pt,a4paper]{amsart}
\usepackage[margin=19mm]{geometry}
\usepackage{amsmath,amssymb,mathtools}
\usepackage[scr=rsfs,cal=euler]{mathalfa}
\usepackage{xfrac}
\usepackage[unicode,hidelinks,bookmarks=false]{hyperref}
\newcommand{\ie}{i.e.\ }
\newcommand{\cf}{cf.\ }
\newcommand{\resp}{respectively\ }
\newcommand{\Subsection}{\subsection}
\providecommand{\nobreakdash}{\nobreak\hbox{-}}
\setlength{\emergencystretch}{2em}
\multlinegap0pt
\usepackage{amscd,mathrsfs}
\usepackage{color,dsfont,pifont,enumitem}
\usepackage{longtable,float,setspace}
\usepackage{multirow,placeins,algorithm,algorithmicx}
\numberwithin{equation}{section}
\newtheorem{lemma}{Lemma}[section]
\newtheorem{defi}{Definition}[section]
\newtheorem{theo}{Theorem}[section]
\newtheorem{coro}{Corollary}[section]
\newtheorem{example}{Example}[section]
\newtheorem{remark}{Remark}[section]
\newtheorem{assumption}{Assumption}[section]
\renewcommand*{\proofname}{\normalfont\bfseries Proof}%
\def\b0{\boldsymbol{0}}
\newcommand{\E}{\mathcal{E}}
\newcommand{\T}{{\mathcal{T}}}
\newcommand{\Q}{{\mathcal{Q}}}
\newcommand{\blambda}{\boldsymbol{\lambda}}
\newcommand{\bPi}{\boldsymbol{\Pi}}
\newcommand{\bPhi}{\boldsymbol{\Phi}}            %new
\newcommand{\bpsi}{\boldsymbol{\psi}}            %new
\newcommand{\ddelta}{\boldsymbol{\delta}}        %new
\newcommand{\sumT}{\sum_{T\in\mathcal{T}_h}}     %new
\newcommand{\sumE}{\sum_{e\in\E_h^0}}            %new
\newcommand{\bw}{\mathbf{w}}
\newcommand{\bu}{\mathbf{u}}
\newcommand{\bv}{\mathbf{v}}
\newcommand{\bl}{\mathbf{L}}
\newcommand{\bh}{\mathbf{H}}
\newcommand{\bi}{\mathbf{I}}
\newcommand{\bn}{\mathbf{n}}
\newcommand{\be}{\mathbf{e}}
\newcommand{\bt}{\mathbf{T}}
\newcommand{\bP}{\mathbf{P}}
\newcommand{\bx}{\mathbf{x}}
\newcommand{\bF}{\mathbf{f}}
\newcommand{\di}{\text{div}}
\newcommand{\eps}{\varepsilon}
\newcommand{\To}{\longrightarrow}
\newcommand{\A}{\mathcal{A}}
\newcommand{\Real}{\mathbb{R}}
\newcommand{\abs}[1]{\left\vert#1\right\vert}
\newcommand{\trb}[1]{|\!|\!|#1|\!|\!|}
\newcommand{\jump}[1]{[\![#1]\!]}
\newcommand{\la}{\langle}
\newcommand{\ra}{\rangle_{\partial T}}
\newcommand{\ran}{\rangle}
\allowdisplaybreaks[4]
\allowdisplaybreaks % For long formula
\setlength{\parindent}{0.25in} \setlength{\parskip}{0.08in}
\begin{document}
\title[The weak Galerkin method for a class of Gross-Pitaevskii typ]{The weak Galerkin method for a class of Gross-Pitaevskii type eigenvalue problems}
\author{Wei Lu; Qilong Zhai}
\date{}
\maketitle
\noindent\emph{Source and licence.} The weak Galerkin method for a class of Gross-Pitaevskii type eigenvalue problems. Version arXiv:2605.29264v1 (CC BY 4.0). This material is distributed under the Creative Commons Attribution 4.0 International licence, http://creativecommons.org/licenses/by/4.0/, as recorded in the arXiv OAI licence metadata for this exact version and shown on the abstract page https://arxiv.org/abs/2605.29264v1. Full rights notice: licenses/arxiv-2605.29264-CC-BY-4.0.txt.\par\smallskip
\noindent\textit{Editorial scope.} Counted unit: the original \texttt{lemma} environment \texttt{lem7} at source lines 334--345 together with its complete proof at lines 346--348; the delivered document keeps source lines 148--349 so that every referenced label and every citation resolves inside it. Native editable source is the paper's own concatenated TeX; the AMS wrapper is editorial formatting only. No publisher class, upstream script or unrelated artwork is redistributed.
\par\smallskip\noindent\textit{Reference note.} This article ships no compiled bibliography (biblatex); the entries delivered here are composed from the author's own BibTeX (.bib) fields, with no field invented and no entry dropped.
\begin{eqnarray}\label{1.1}
	\inf\{E(v):v\in H_0^1(\Omega),\displaystyle\int_\Omega v^2 d\Omega=1\},
\end{eqnarray}
where $\Omega\subset\Real^2$ is a convex bounded domain, and the energy function $E$ is defined as
\begin{eqnarray}\label{E}
	E(v)=\frac12 a(v,v)+\frac12\displaystyle\int_\Omega F(v^2)d\Omega,
\end{eqnarray}
where
\begin{eqnarray*}
	a(u,v)=(\A\nabla u,\nabla v)+(Vu,v),\quad\forall u,v\in H^1(\Omega),
\end{eqnarray*}
with $\A(x)\in (L^\infty(\Omega))^{2\times2}$ is symmetric, and 
%$V(x)\in L^p(\Omega)(1<p\le\infty)$
$V(x)\in L^2(\Omega)$ is the non-negative potential.

Next, we introduce the following assumptions.

\textbf{Assumption (A1).}\emph{There exists a constant $\alpha$ such that $\displaystyle\inf_{x\in \Omega}\xi^T\A(x)\xi\ge\alpha\xi^T\xi$, $\forall \xi\in R^2$, $\forall x\in\Omega$ a.e..
}

\textbf{Assumption (A2).}\emph{$F(t)\in C^1[0,+\infty)\cap C^2(0,+\infty)$, and $\displaystyle\inf_{t>0}F''(t)>0$.
}

\textbf{Assumption (A3).}\emph{There exists $q\in [0,2)$ and a constant $C$ such that $|F'(t)|\le C(1+t^q)$, $\forall t\ge0$.
}

\textbf{Assumption (A4).}\emph{$F''(t)t$ locally bounded in $[0,+\infty)$.
}

\textbf{Assumption (A5).}\emph{There exists $1<r\le 2$ and $0\le s\le 3-r$ such that: $\forall R>0$, there holds $\forall t_1\in(0,R)$ and $\forall t_2\in\Real$ that
	\begin{eqnarray*}
		|F'(t_2^2)t_2-F'(t_1^2)t_2-2F''(t_1^2)t^2(t_2-t_1)|\le C(1+|t_2|^s)|t_2-t_1|^r,
	\end{eqnarray*}
	where C is a positive constant depend on R.
}


For simplicity, we denote by $f(t)=F'(t)$. Then, under Assumption (A2), $E$ is twice differentiable on $H^1(\Omega)$ and we have
\begin{eqnarray*}
	E'(v)=A^vv,\quad \la E''(u)v,w\ran=\la A^uv,w\ran+2(f'(u^2)u^2v,w),\quad\forall u,v,w\in H^1(\Omega),
\end{eqnarray*}
where
\begin{eqnarray*}
	A^vw=-\di(\A\nabla\cdot w)+Vw+f(v^2)w,\quad\forall v,w\in H^1(\Omega).
\end{eqnarray*}
It is easy to check, for each $u\in L^2(\Omega)$, $A^u$ defines a self-adjoint operator and satisfies
\begin{eqnarray*}
	\la A^uv,w\ran=(\A\nabla v,\nabla w)+(Vv,w)+(f(u^2)v,w),\quad\forall v,w\in H_0^1(\Omega).
\end{eqnarray*}

\begin{lemma}\label{lem1}
	Under Assumptions (A1)-(A3), \eqref{1.1} has a unique (up to sign) minimizer $u$. Moreover, there exists $\lambda\in\Real$ such that $(u,\lambda)$ is the first eigenpair of the following nonlinear eigenvalue problem:
	\begin{equation}\label{1.3}
		\left\{\begin{array}{rcl}
			\la A^uu,v\ran&=&\lambda (u,v),\quad\forall v\in H_0^1(\Omega),\\
			\|u\| &=&1.
		\end{array}\right.
	\end{equation}
	Besides, $u\in C(\bar{\Omega})\cap L^\infty(\Omega)$ is positive in $\Omega$, and $\lambda$ is a simple eigenvalue of $A^u$.
\end{lemma}
\begin{proof}
	See Lemma 2 in \cite{Eric2009}.
\end{proof}

\section{Weak Galerkin discretization}In this section, we introduce the standard WG scheme for nonlinear eigenvalue problem \eqref{1.3}.


Let $\T_h$ be a partition of the domain $\Omega$, and the elements
in $\T_h$ are polygons satisfying the regular assumptions specified
in \cite{Wang2014a}. Let $\E_h$ be the edges in $\T_h$, and
$\E_h^0$ denotes by the interior edges $\E_h \backslash\partial\Omega$. For each element $T\in\T_h$, $h_T$ represents the diameter of $T$, and $h=\max\limits_{T\in\T_h} h_T$ denotes the mesh size. For simplicity, we use $a\lesssim b$  and $a\gtrsim b$ to represent $a\le Cb$ and $a\ge Cb$, respectively, where $C$ is a constant independent of mesh size $h$.

%According to \cite{Wang2013a}, the following trace inequality and inverse inequality hold true.
\begin{lemma}
	For any element  $T\in\T_h$, there holds the trace inequality
	\begin{eqnarray}\label{trace}
		\|v\|_{\partial T}^2\lesssim h_T^{-1}\|v\|_T^2+h_T\|\nabla v\|_T^2,\quad\forall v\in h^1(T).
	\end{eqnarray}
\end{lemma}
\begin{proof}
	See Lemma A.3 in \cite{Wang2014a}.
\end{proof}

%\begin{lemma}
%	For any element $T\in\T_h$, there holds the inverse inequality
%	\begin{eqnarray}\label{inverse}
	%		\|\nabla v\|_T\lesssim h_T^{-1}\|v\|_T,\quad\forall v\in P_k(T).
	%	\end{eqnarray}
%\end{lemma}

Then, we introduce the WG discretization for the eigenvalue problem \eqref{1.3}. For a given integer $k\ge1$, the WG finite element space is defined by
\begin{eqnarray*}
	V_h= \big\{ v=\{ v_0, v_b\}: v_0|_T\in  P_k(T),  v_b|_e\in  P_{k-1}(e),\ \forall T\in\T_h,\ e\in\E_h,\text{ and }  v_b|_{\partial\Omega}= 0\big\},
\end{eqnarray*}

%Define the sum space $V=V_h+ H_0^1(\Omega)$. For each $ v\in V$, we define its weak gradient $\nabla_w v$ as follows.
% \begin{defi}
	%	$\nabla_w v|_T$ is the unique polynomial in $[P_{k-1}(T)]^{2}$ satisfying
	%	 \begin{eqnarray*}
		%		(\nabla_w v,q)_T=-(v_0,\nabla\cdot q)_T+\la v_b,q\cdot\n
		%		\ra,\quad\forall q\in [P_{k-1}(T)]^{2},
		%	\end{eqnarray*}
	%	where $\n$ denotes the outward unit normal vector.
	%\end{defi}
	
	For each $ v_h\in V_h$, we define its weak gradient $\nabla_w v_h$ as follows.
	\begin{defi}
		For each $ v_h\in V_h$, $\nabla_w v_h|_T$ is the unique polynomial in $[P_{k-1}(T)]^{2}$
		satisfying
		\begin{eqnarray*}
			(\nabla_w v_h,q)_T=-(v_0,\nabla\cdot q)_T+\la v_b,q\cdot\bn
			\ra,\quad\forall q\in [P_{k-1}(T)]^2,
		\end{eqnarray*}
		where $\bn$ denotes the outward unit normal vector.
	\end{defi}
	
	For each $T\in\T_h$, let $Q_0$ denotes the $L^2$ projection from $L^2(T)$ onto $ P_k(T)$, $\Q_h$ denotes the $L^2$ projection from $[L^2(T)]^{2}$ onto $[P_{k-1}(T)]^{2}$.
	For each $e\in\E_h$, let $Q_b$ denotes the $L^2$ projection from $L^2(e)$ onto $ P_{k-1}(e)$ for each $e\in\E_h$.
	Combining $Q_0$ and $Q_b$ together, we define $Q_h=\{Q_0,Q_b\}$,
	which is a projection onto $V_h$.
	
	For the projection operators defined above, we have the following results.
	\begin{lemma}\label{lem3}
		There holds the following commutative property on each $T\in\T_h$:
		\begin{eqnarray*}
			\nabla_wQ_hv=\Q_h\nabla v, \quad\forall v\in H^1(T).
		\end{eqnarray*}
	\end{lemma}
	\begin{proof}
		See Lemma 5.1 in \cite{Mu2013c}.
	\end{proof}
	
	%Refer to \cite{Mu2013c}, the following estimates hold true for the projection operators.
	\begin{lemma}
		For any $v\in H^{k+1}(\Omega)$ and $\tau\in [H^k(\Omega)]^2$, $s=0,1$, we have
		\begin{eqnarray}
			\sumT h^{2s}\|v-Q_hv\|_{s,T}^2&\lesssim h^{2(k+1)}\|v\|_{k+1}^2,\label{pro1}\\
			\sumT h^{2s}\|\tau-\Q_h\tau\|_{s,T}^2&\lesssim h^{2k}\|\tau\|_k^2.\label{pro2}
		\end{eqnarray}
	\end{lemma}
	\begin{proof}
		See Lemma 4.1 in \cite{Wang2014a}.
	\end{proof}
	
	In order to give the WG scheme, we define two bilinear forms on $V_h$.
	\begin{eqnarray}
		s( v_h, w_h)=&\sumT h_T^{-1+\eps}\langle  Q_bv_0- v_b, Q_bw_0- w_b\rangle_{\partial T},\label{2.1}\\
		a_h( v_h, w_h)=&(\A\nabla_wv_h,\nabla_ww_h)+(Vv_h,w_h)+s(v_h,w_h),\label{2.2}
	\end{eqnarray}
	where $v_h=\{v_0,v_b\}$, $w_h=\{w_0,w_b\}\in V_h$, $\eps\in(0,1)$ is a small constant.
	
	In addition, for each $u_h\in V_h$, we introduce the operator $A_h^{u_h}$ by
	\begin{eqnarray}\label{2.3}
		\la A_h^{u_h}v_h, w_h\ran=a_h( v_h, w_h)+(f(u_h^2)v_h,w_h), \quad\forall v_h,w_h\in V_h,
	\end{eqnarray}
	and the discrete energy function is given by
	\begin{eqnarray}\label{Eh}
		E_h(v_h)=\frac12 a_h(v_h,v_h)+\frac12\displaystyle\int_\Omega F(v_h^2)d\Omega, \quad\forall v_h\in V_h.
	\end{eqnarray}
	Similarly, $E_h$ is twice differentiable on $V_h$ and satisfies
	\begin{eqnarray}\label{E''h}
		\la E_h''(u_h)v_h,w_h\ran=\la A_h^{u_h}v_h,w_h\ran+2(f'(u_h^2)u_h^2v_h,w_h),\quad\forall u_h, v_h,w_h\in V_h.
	\end{eqnarray}
	
	Now, we are ready to give the weak Galerkin scheme for the nonlinear eigenvalue problem \eqref{1.3}.
	\begin{algorithm}[ht!]
		\caption{The weak Galerkin scheme for \eqref{1.3}.}
		Find $ u_h\in V_h$ and $\lambda_h\in\Real$ such that $\| u_h\|=1$ and
		\begin{eqnarray}\label{2.4}
			\la A_h^{u_h} u_h, v_h\ran=\lambda_h(u_h,v_h),\quad\forall  v_h\in V_h.
		\end{eqnarray}
	\end{algorithm}

\section{Error analysis}In this section, we show the error analysis for the WG scheme \eqref{2.4}.

For the aim of analysis, we introduce
\begin{eqnarray*}
	\|v_h\|_{1,h}^2=\sumT\|\nabla v_0\|_T^2+\sumT h_T^{-1}\|  Q_bv_0- v_b\|_{\partial T}^2,\quad\forall v_h\in V_h.
\end{eqnarray*}
It is easy to check $\|\cdot\|_{1,h}$ defines a norm on $V_h$ \cite{Zhai2019c,Wang2017b}.

Consider the following auxiliary problem: Find $\widetilde{u}_h\in V_h$ and $\widetilde{\lambda}_h\in\Real$ such that
\begin{eqnarray}\label{3.3}
	\la A_h^{u} \widetilde{u}_h, v_h\ran=\widetilde{\lambda}_h(\widetilde{u}_h,v_h),\quad\forall  v_h\in V_h,
\end{eqnarray}
where $u$ is the non-negative minimizer of \eqref{1.1}. Then, we have the following estimates.

\begin{lemma}\label{lem7}
	Let $(\widetilde{u}_h,\widetilde{\lambda}_h)$ be the first eigenpair of \eqref{3.3}. Suppose the ground state $u\in H^{k+1}(\Omega)$, then we have
	\begin{eqnarray}
		h^{2k}\|u\|_{k+1}\lesssim\lambda-\widetilde{\lambda}_h&\lesssim h^{2k-2\eps}\|u\|_{k+1},\label{3.4}
		%		|\lambda-\widetilde{\lambda}_h|&\lesssim h^{2k-2\eps}\|u\|_{k+1},\label{3.4}
		\\
		\|u-\widetilde{u}_h\|_{1,h}&\lesssim h^{k-\eps}\|u\|_{k+1},\label{3.5}
		\\
		\|u-\widetilde{u}_h\|&\lesssim h^{k+1-\eps}\|u\|_{k+1},\label{3.6}
	\end{eqnarray}
	when the mesh size h is sufficiently small.
\end{lemma}

\begin{proof}
	See \cite{Zhai2019c}.
\end{proof}


\begin{thebibliography}{99}
\bibitem[]{Eric2009} Eric Cancès; Chakir, Rachida; Maday, Yvon. Numerical Analysis of Nonlinear Eigenvalue Problems. Journal of Scientific Computing 45 (2009)
\bibitem[]{Mu2013c} Mu, Lin; Wang, Junping; Ye, Xiu. A weak Galerkin finite element method with polynomial reduction. J. Comput. Appl. Math. 285 (2015) 45--58
\bibitem[]{Wang2014a} Wang, Junping; Ye, Xiu. A weak Galerkin mixed finite element method for second order elliptic problems. Math. Comput. 83 (2014) 2101--2126
\bibitem[]{Wang2017b} Wang, J.; Wang, R.; Zhai, Q.; Zhang, R.. A Systematic Study on Weak Galerkin Finite Element Methods for Second Order Elliptic Problems. J. Sci. Comput. (2017)
\bibitem[]{Zhai2019c} Zhai, Qilong; Xie, Hehu; Zhang, Ran; Zhang, Zhimin. The weak Galerkin method for elliptic eigenvalue problems. Communications in Computational Physics 26 (2019) 160--191
\end{thebibliography}
\end{document}
